% chapters/ch32_u_def_01.tex
% Chapter 32 (Part IV): U-DEF-01 Addendum — Conflict Resolution
% Source: NRMO_v5_updated.pdf, Part IV (lines 10302-10334 of v5_layout.txt)

\chapter{U-DEF-01 Addendum --- Conflict Resolution}
\label{ch:u-def-01}

\nrmobox{Document identification}{
\textbf{Title}: NRMO Unified Specification | Addendum U-DEF-01\\[0.3em]
\textbf{Subtitle}: NRMO Unified Specification Theoretical Unification
Addendum\\[0.3em]
\textbf{Insertion}: Between Part~I (v2.1) and Part~III
(v4 Theoretical Foundation)\\[0.3em]
\textbf{Author}: Takashi Ikeya | \textbf{Conflict ID}: U-DEF-01 |
\textbf{Status}: ADOPTED | \textbf{Year}: 2026
}

% =====================================================================
\section{Identified Conflict: U-DEF-01}
\label{sec:u-def-01-conflict}

A single formal conflict exists between the two constituent
documents of this specification. All other design elements ---
Phased Parallel Execution, $\varepsilon$-buffer update, Governance
Toolkit, APSOC, DEADLOCK handler --- are mutually consistent.

\begin{table}[ht]
\centering
\small
\begin{tabularx}{\linewidth}{lXX}
\toprule
\textbf{Dimension} & \textbf{Part~I v2.1 (p.\,19)} & \textbf{Part~III / v4 Paper Def.\,3} \\
\midrule

$U(\pi)$ definition
& ``Evaluated on survival'' --- conditional on $\tau_F = \infty$
& $\liminf (1/T) E_\pi[\sum_{t=0}^{T-1} r_t]$ --- unconditional \\

Risk
& Zero-probability conditioning when $\mathrm{ruin} \to 1$
& None; well-defined for all $\pi$ \\

Role
& Primary optimisation target (implicit)
& Primary optimisation target (explicit) \\

\bottomrule
\end{tabularx}
\caption[U-DEF-01 conflict]{The U-DEF-01 conflict: implicit vs.~
explicit optimisation target with different domain conditioning.}
\label{tab:u-def-01-conflict}
\end{table}

% =====================================================================
\section{Resolution: Plan A}
\label{sec:u-def-01-resolution}

This Addendum adopts Plan A: \textbf{unconditional $U(\pi)$ is the
canonical optimisation scalar}. The conditional form
$U_{\mathrm{cond}}(\pi)$ is retained as an audit metric only.

\paragraph{Decree U-DEF-01-A.}

\begin{equation}
U(\pi) := \liminf_{T \to \infty}
\frac{1}{T} \mathbb{E}_\pi\!\left[
\sum_{t=0}^{T-1} r(s_t, a_t)
\right]
\quad \text{(unconditional)}.
\label{eq:u-def-01-canonical}
\end{equation}

This supersedes Part~I p.~19 (``evaluated on survival'').

\paragraph{Rationale.}

\begin{enumerate}
\item Well-defined for all $\pi$ including high-ruin policies.
\item Directly aligns with Theorem~1 (Ruin Dilution): $U(\pi) \le
(1 - q) R_{\max}$ makes ruin-cost visible in the objective itself.
\item For $\pi \in \mathrm{safe}(T, 0)$, unconditional $=$
conditional since $P(\tau_F = \infty) = 1$.
\end{enumerate}

\paragraph{Prohibition U-DEF-01-P.}
Using $U_{\mathrm{cond}}(\pi)$ as optimisation target in Algorithm~1
Step~9 or as APSOC gate input is \textbf{prohibited}
(equivalent to Chapter~33 Prohibition~\#1).

% =====================================================================
\section{Role Assignment After Resolution}
\label{sec:u-def-01-role-assignment}

\begin{table}[ht]
\centering
\small
\begin{tabularx}{\linewidth}{lXXc}
\toprule
\textbf{Metric} & \textbf{Role} & \textbf{Computed in} & \textbf{Optimisation?} \\
\midrule

$U(\pi)$ unconditional
& Primary scalar
& Algorithm~1 Step~9
& YES \\

$U_{\mathrm{cond}}(\pi)$ conditional
& Audit / interpretation metric
& Evidence Capsule (read-only)
& NO \\

$q = P(\tau_F \le T)$
& Constraint certificate
& APSOC Gate / CVaR
& Constraint only \\

\bottomrule
\end{tabularx}
\caption[U-DEF-01 role assignment]{Role assignment after U-DEF-01-A
resolution. Only $U(\pi)$ is an optimisation target.}
\label{tab:u-def-01-roles}
\end{table}

% =====================================================================
\section*{Connections to Other Parts}

\begin{itemize}
\item U-DEF-01-A directly affects Part~V Theorem~1 (Ruin Dilution):
the unconditional definition makes the bound $U(\pi) \le (1 - q)
R_{\max}$ a direct theorem, not a conditional one.
\item Algorithm~1 (Phased Parallel Execution) in Part~V Step~9
returns $\arg\max_{o \in O_{\mathrm{valid}}} U(\pi_o)$ ---
unconditional.
\item APSOC Gate input (CVaR ruin upper bound) uses the constraint
side $q$, not $U_{\mathrm{cond}}$.
\item In the engineering implementation (Part~X), the v2.5 codebase
metrics module computes both $U$ and $U_{\mathrm{cond}}$;
$U_{\mathrm{cond}}$ is logged audit-only, $U$ feeds the scoring
loop.
\end{itemize}
